% ECONOMICS_PROBLEM: {"id": "D61-402", "title": "Welfare effects of nudges", "jel_code": "D61", "source_id": "OP-0309"}
% !TeX program = xelatex
\documentclass[11pt,a4paper]{article}
\usepackage[margin=23mm]{geometry}
\usepackage{fontspec}
\setmainfont{Latin Modern Roman}
\usepackage{amsmath,amssymb,mathtools}
\usepackage{xcolor,enumitem,booktabs,longtable,array,xurl,hyperref}
\definecolor{muted}{HTML}{53636C}
\hypersetup{colorlinks=true,urlcolor=blue,linkcolor=blue}
\setlength{\parindent}{0pt}
\setlength{\parskip}{5pt}
\newcommand{\R}{\mathbb R}
\newcommand{\E}{\mathbb E}
\newcommand{\N}{\mathbb N}
\newcommand{\ind}{\mathbf 1}
\newcommand{\Var}{\operatorname{Var}}
\newcommand{\Cov}{\operatorname{Cov}}
\newcommand{\supp}{\operatorname{supp}}
\DeclareMathOperator*{\argmin}{arg\,min}
\DeclareMathOperator*{\argmax}{arg\,max}
\newcommand{\entrymeta}[3]{{\sffamily\small\color{muted}\textbf{\texttt{#1}}\quad|\quad #2\par #3\par}}
\newcommand{\Problemref}[1]{\texttt{#1}}
\newcommand{\JELref}[2]{\texttt{#2} (JEL \texttt{#1})}
\begin{document}
\section*{Welfare effects of nudges}
\textbf{Problem ID:} \texttt{D61-402}\quad \textbf{JEL:} \texttt{D61}\par
% BEGIN ECONOMICS STATEMENT
\label{problem:OP-0309}
{\sffamily\small\color{muted}Original subject group: Behavioral and experimental economics\par}
\label{definitions:problem:30007591}
\textbf{Audit classification:} Background; proposed extension. Allcott et al. 2025 is genuine and its abstract presents welfare conditions as results. The general application-level identification specification is editorial; no remaining exact theorem was verified.
\subsection*{Definitions and precise research target}
\textbf{Model and research target.} Consider consumers \(i\) choosing quantities \(q_i\) of a good. Let \(v_i(q)\) be experienced monetary-equivalent benefit, \(b_i(q,n)\) a choice distortion under nudge \(n\in\{0,1\}\), and \(p(n,t)\) the equilibrium price when a tax \(t\) is fixed. A specified behavioral demand model is
\[q_i(n,t)\in\arg\max_{q\geq0}\{v_i(q)+b_i(q,n)-(p(n,t)+t)q\}.\]
Define welfare, including production costs \(C\), aggregate external damage \(H\), and real implementation cost \(K\), by
\[W(n,t)=\mathbb E[v_i(q_i(n,t))]-C(Q(n,t))-H(Q(n,t))-K(n),\qquad Q(n,t)=\mathbb E[q_i(n,t)].\]
Transfers from prices and tax receipts cancel only under the chosen welfare weights and incidence assumptions. The target estimand is \(\Delta W(t)=W(1,t)-W(0,t)\), or the corresponding difference with taxes reoptimized under each nudge. Determine which experiments and market data identify experienced benefits, heterogeneous distortions, price responses, and spillovers sufficiently to sign \(\Delta W\) in a specified policy application. A mean change in purchases does not by itself identify welfare, and observed choice utility need not equal experienced benefit. Report sensitivity to benefit measurement, welfare weights, unobserved heterogeneity, and equilibrium assumptions.

\emph{Definitions and maintained assumptions (editorial unless a source definition is named).} Normalize consumers to mass one on a probability space, with \(q_i\in[0,\bar q]\), finite \(\bar q\), integrable experienced benefits and continuous choice payoff. Fix a measurable tie-breaking rule for demand if maximizers are multiple. Supply has a specified competitive cost function \(C\), market clearing and an equilibrium-selection rule determining \(p(n,t)\); price cannot merely be named an equilibrium without that supply model. \(v_i\) is the normative money-equivalent benefit and \(b_i\) a decision distortion, not a welfare benefit unless the declared criterion includes it. \(C(Q),H(Q),K(n)\) are total real costs per unit consumer mass, with \(H\) external damage. Equal welfare weights and public revenue rebated as transfers justify dropping prices and taxes from total surplus; unequal weights require separately counting incidence. Fixed-tax \(\Delta W(t)\) differs from \(\sup_{t\in\mathcal T}W(1,t)-\sup_{t\in\mathcal T}W(0,t)\), where tax set \(\mathcal T\) and fiscal constraints are declared. Observed purchase changes do not identify \(v_i\), distortions or cost functions without restrictions or auxiliary measurements.
\subsection*{Variants and assumption changes}
\begin{enumerate}
\item \textbf{Fixed tax/price experiment (editorial).} Randomize labels holding product prices and taxes fixed, measure experienced benefits by a defended protocol, and bound welfare sensitivity to distortion heterogeneity and implementation costs.
\item \textbf{Equilibrium and tax optimization (editorial).} Specify supply elasticity, pass-through and a compact admissible tax set. Identify nudge welfare after equilibrium response and compare separately with taxes optimized under each regime, preserving the same public budget and welfare weights.
\end{enumerate}
\subsection*{References and source audit}
\begin{enumerate}
\item Publisher abstract states welfare results; full PDF inaccessible. No exact remaining identification question verified. \url{https://www.aeaweb.org/articles?id=10.1257/aer.20231304}
\end{enumerate}
\begin{enumerate}\raggedright
\item\hypertarget{definitions:source:30007591:1}{} Hunt Allcott; Daniel Cohen; William Morrison; Dmitry Taubinsky. 2025. \emph{When Do "Nudges" Increase Welfare?}. Abstract; AER 115(5), pp. 1555–1596. Full paper was not accessible in this audit..
\url{https://www.aeaweb.org/articles?id=10.1257/aer.20231304}.
DOI: \url{https://doi.org/10.1257/aer.20231304}.
\end{enumerate}

\subsection*{Common conventions: Source mathematical conventions}

These conventions apply to each entry unless explicitly replaced.

\textbf{Models and identification.} A primitive model \(m\) specifies all probability laws, feasible choices, information, equilibrium and measurement rules used by its target. The admissible class \(\mathcal M\) is fixed independently of outcomes. If \(P_O(m)\) is its observable probability law and \(\tau(m)\) the target, its identified set at observable law \(P\) is
\[
 \mathcal I_\tau(P)=\{\tau(m):m\in\mathcal M,\ P_O(m)=P\}.
\]
Point identification means that this set is a singleton. An empty set signals incompatibility of data and assumptions. Coordinate bounds need not describe the joint set. Bounds are sharp only if every retained target is attainable or arbitrarily approachable by compatible models. Finite bounds require explicit restrictions on unobserved quantities; unrestricted confounding, hidden wealth, extrapolation or arbitrary equilibrium selection can yield uninformative sets.

\textbf{Causal interventions.} A potential outcome \(Y_i(p)\) is the outcome under a complete policy regime \(p\), including timing, financing, information and market spillovers. The target population and units are fixed before treatment. With interference, use the complete assignment vector or a justified exposure mapping; individual no-interference notation alone cannot identify an economy-wide intervention. A difference of mean potential outcomes does not require the cross-world joint distribution. Conditioning on post-treatment migration, survival, enrollment or employment changes the estimand and needs separate justification.

\textbf{Statistical statements.} An estimator, confidence set, forecast or test specifies a sampling law, model class and loss/coverage criterion. Uniform \((1-\alpha)\)-coverage means \(\inf_{m\in\mathcal M}\Pr_m\{\tau(m)\in C_n\}\geq1-\alpha\), or an explicitly asymptotic version. The whole parameter space trivially has coverage, so informativeness requires a stated width, risk or power objective. A forecasting improvement is relative to a named benchmark, information set and evaluation population.

\textbf{Equilibrium and optimization.} An equilibrium comprises feasible plans, optimality, beliefs where needed, budget constraints and all market/resource clearing. The primitives or a stated benchmark define each correspondence \(\mathcal E(p;m)\). Nonemptiness is assumed or proved, never inferred just from compact policy sets. A selected-equilibrium target uses a declared selection rule; a robust target quantifies over all permitted equilibria. Use suprema or infima unless attainment is established. An \(\varepsilon\)-optimal policy is feasible and within \(\varepsilon\) of the finite optimum value. Report an empty feasible set separately.

\textbf{Welfare and accounting.} Normative utility, interpersonal normalization, social weights, numeraire, discounting and the use of public receipts are primitives. Behavioral utility need not coincide with the welfare criterion. Household welfare includes ownership income and rents once; adding profits again to compensating variation generally double-counts them. A monetary equivalent is defined by an explicit transfer and price convention, with monotonicity and range conditions. Average effects, mechanism contributions, transfers and welfare losses are different targets. Infinite sums require convergence and dynamic feasible plans require terminal or no-Ponzi conditions.

\textbf{Source versions.} A working-paper date, revision date and journal publication year are distinguished. A DOI for a journal version is not automatically the DOI of an earlier manuscript. Locators refer to the stated version and printed pages where specified. A metadata-only check does not verify a claim about a full-text conclusion. Every entry's source subsection narrows the provenance claim accordingly.
% END ECONOMICS STATEMENT
\end{document}
